\subsection{Updated Constraints on Large Extra Dimensions from Reactor Antineutrino Experiments --- arXiv:2510.12900}

\textbf{Authors:} T. G\"okalp Ela\c{c}maz, Ivan Martinez-Soler, Yuber F. Perez-Gonzalez

\paragraph{主な主張}
Daya Bay, RENO, KamLAND, NEOS, STEREO の原子炉反ニュートリノデータから、最軽量 active neutrino を質量ゼロとした $d=1$ large-extra-dimension 模型に対して $a\lesssim0.58\,\mu\mathrm{m}$（normal ordering）および $a\lesssim0.12\,\mu\mathrm{m}$（inverted ordering）の 99\% C.L. 制限を得るとともに、等半径の $d=2,3,4$ では余剰次元数の増加とともに制限が強くなることを示した \cite{Elacmaz:2025LED}。

\paragraph{新規性}
最新の複数 reactor data set を統合して $d=1$ の従来制限を更新し、さらに等半径を仮定した $d=2,3,4$ の KK tower を同じ枠組みで系統的に解析して新たな制限を与えた点が新しい。

\paragraph{位置付け}
bulk right-handed neutrino を持つ large-extra-dimension 模型を低エネルギーの $\bar\nu_e$ disappearance から検証する研究であり、一般の torus geometry や complex-structure dependence を考える際の等半径 compactification に対する実験的 baseline と位置付けられる。

\paragraph{コメント：torus geometry の仮定}
$d=2$ では $m_{n_1,n_2}^2=(n_1^2+n_2^2)/a^2$ の KK spectrum を用いており、これは等半径かつ直交する square torus、すなわち $\tau=i$ に対応する特殊な場合である。したがって一般の complex structure modulus $\tau$ を動かした解析は行っておらず、shape dependence を含めるには $m_{mn}^2\propto |m-n\tau|^2/\tau_2$ のような一般の $T^2$ KK spectrum への拡張が必要である。
